\documentclass[
    paper=a4,
    fontsize=11pt,
    twoside=false
]{scrreprt}

\usepackage[T1]{fontenc}
\usepackage[utf8]{inputenc}
\usepackage{lmodern}
\usepackage[brazilian]{babel}
\usepackage[margin=2.5cm]{geometry}
\usepackage{microtype}
\usepackage[dvipsnames]{xcolor}
\usepackage{hyperref}

\hypersetup{
    colorlinks=true,
    linkcolor=RoyalBlue,
    urlcolor=RoyalBlue,
    pdftitle={Relatório Cloud LaTeX},
    pdfauthor={José Dercy Silva}
}

\begin{document}

\title{Relatório Cloud LaTeX}
\subtitle{Ambiente de Alta Performance com Tectonic e GitHub Actions}
\author{José Dercy Silva}
\date{\today}
\maketitle

\chapter{Introdução}
Este documento atesta a configuração bem-sucedida do repositório local. O código-fonte está sendo versionado via Git e a arquitetura está preparada para integração contínua.

\section{Próximos Passos}
A próxima fase da infraestrutura envolverá a automação da compilação utilizando os servidores da Microsoft via GitHub Actions.

\end{document}
